\documentclass[11pt]{article}
\usepackage[margin=0.82in]{geometry}
\usepackage[T1]{fontenc}
\usepackage{lmodern,amsmath,amssymb,amsthm,booktabs,hyperref,array}
\hypersetup{colorlinks=true,linkcolor=blue,urlcolor=blue}
\newtheorem{theorem}{Theorem}
\newtheorem{lemma}{Lemma}
\newtheorem{proposition}{Proposition}
\newcommand{\C}{\mathrm C}
\newcommand{\J}{\mathrm J}
\newcommand{\Q}{\mathrm Q}
\newcommand{\K}{\mathrm K}
\renewcommand{\AA}{\mathrm A}
\setlength{\parindent}{0pt}
\setlength{\parskip}{5pt}
\emergencystretch=2em
\title{Long Periodic Games of Beggar-My-Neighbour:\\Balanced Deals and Protected Queue Modules}
\author{Christopher D. Long}
\date{19 November 2026}
\begin{document}
\maketitle
\begin{abstract}
We give explicit periodic games for the standard $52$-card inventory of
Beggar-My-Neighbour, under chronological collection and winner-led play.
A $26$--$26$ deal enters a primitive rank-state cycle of $156$ completed
tricks and $1{,}284$ individual plays after two tricks and $22$ plays.
An unrestricted $38$--$14$ starting state lies directly on a $200$-trick,
$1{,}264$-play cycle; a related $50$--$2$ construction has period $176$
tricks and $1{,}228$ plays. The balanced construction is organized by a
protected queue-module operation: one $12$-phase module and two fixed
spacers give a $13$-trick control round repeated twelve times.
We give a general two-trick balanced-loader lemma and describe exhaustive
backward-component certificates for all three orbits. All existence and
period assertions are accompanied by a standalone, dependency-free verifier.
No global maximality of these periods is claimed.
\end{abstract}

\section{Rules and notions of period}
A hand is a first-in, first-out queue, written with its next card at the
left. An ordinary card is denoted $\C$; $\J,\Q,\K,\AA$ require respectively
$1,2,3,4$ payments by the opponent. A face card played during a payment
interrupts that payment and becomes the new claim. A completed payment is
collected in chronological play order by the claimant, who leads the next
trick. These are the conventions used in \cite{casella} and in the
companion minimum-period work \cite{longshort}.

We consider the inventory
\[
  \C^{36}\J^4\Q^4\K^4\AA^4.
\]
A \emph{rank state} is $(H_A,H_B,\ell)$ at an empty-table trick boundary,
where $\ell\in\{A,B\}$ is the next leader. Ordinary cards and cards of
the same face rank are indistinguishable; the two players are named.
The \emph{primitive trick period} of a recurrent state is the smallest
$p>0$ for which the entire rank state returns after $p$ completed tricks.
Its \emph{play period} counts individual cards played during that return.
An initial transient is excluded from both periods. We distinguish cycles
accessible from a \emph{balanced} $26$--$26$ deal from cycles requiring an
unrestricted starting split. A return of the rank state need not be a
return of all $52$ individually labelled cards.

Casella et al. constructed a cycle of $62$ tricks and $440$ plays accessible
from a balanced deal \cite{casella}; further dynamical investigations are
reported in \cite{andorno}. Our results give explicit longer examples, not
claims of optimality among all decks.

\section{Three explicit long cycles}
\begin{theorem}[Balanced cycle]\label{thm:balanced}
With $A$ leading, the $26$--$26$ deal
\begin{verbatim}
A: CCCCCCCCCCJCCCCCCCCQCKCKQK
B: CCCCCCCCJCCKQCAAQAAJCCJCCC
\end{verbatim}
reaches a cycle after exactly two completed tricks, containing $19$ and
$3$ plays respectively. The cycle has primitive period
\[
\boxed{156\ \text{tricks},\qquad 1{,}284\ \text{plays}.}
\]
\end{theorem}
\begin{theorem}[Unrestricted cycles]\label{thm:other}
The following states have the full standard inventory and lie directly
on cycles, with $A$ leading in each case:
\begin{center}
\begin{tabular}{@{}lll@{}}
\toprule
Initial split & Primitive tricks & Plays per period\\
\midrule
$38$--$14$ & $200$ & $1{,}264$\\
$50$--$2$ & $176$ & $1{,}228$\\
\bottomrule
\end{tabular}
\end{center}
The exact queues are
\begin{verbatim}
200 tricks:
A: CCCCAQCACKJCCACCCCCJCCCCCACCJCCKCCCQCC
B: CCQCKCCCQCKCJC

176 tricks:
A: CCKCCCQCCCCCCCCCCKCKQCQCKQCACAJCCACACCCJCCCCCCCCJC
B: JC
\end{verbatim}
\end{theorem}

The inventory and primitive periods in these theorems can be checked by
ordinary forward simulation. The supplied program \texttt{verify.py}
uses two structurally different simulations---a card-index implementation
and a challenge-block implementation---and compares their results at every
trick boundary. It also confirms that every boundary state within each
reported period is distinct. This finite check is a reproducible proof of
the two assertions, independent of the discovery procedure.

\section{The protected $\J\C$ module}
The balanced example has a more explanatory word certificate.

\begin{lemma}[Protected concatenation]\label{lem:protected}
Let $P$ contain no jacks. Start with the large queue $P\J*$ and the
reservoir queue $\J\C$, the large player leading, where $*$ is a
protected marker. Suppose ordinary play finishes a sequence of completed
tricks without reading $*$ and reaches a boundary at which the reservoir
again equals $\J\C$, the large player leads, and their queue is
$*X\J\C$. Then the \emph{successor module} is $\C X$.
These passages may be concatenated: replacing the protected marker by
subsequent modules preserves the completed passage and its cost.
\end{lemma}
\begin{proof}
Before reaching the protected marker, no card to its right can affect any
payment or choice of next player. The trick sequence is therefore independent
of the suffix hidden behind $*$. When the marker reaches the front, the
remaining queue is $X\J\C$ behind it and the reservoir and leader have
returned. In the concatenated deck, the protected suffix begins with the
ordinary card $\C$ separating adjacent modules, giving the next piece
$\C X$. Replacing $*$ by this suffix changes no earlier play.
\end{proof}

For Theorem~\ref{thm:balanced}, use the module and two fixed spacers
\begin{verbatim}
P = CCKCCCQCCCCQCKKCKQCAAQAA
U = CC
V = CCCCCCCCCCCCCCCCCCCC
\end{verbatim}
The module inventory is $\C^{12}\Q^4\K^4\AA^4$. At a convenient cycle
phase, the small queue is $\J\C$, the large player leads, and the large
queue is
\[
P\J\,U\J\,V\J\C.
\]
There are twelve distinct successive modules $P=P_0,P_1,\ldots,P_{11}$
under Lemma~\ref{lem:protected}, returning to $P_{12}=P_0$. Each passage
of $P_j$ takes exactly seven tricks and $59$ plays, and has the same
winner word. The two spacers are fixed points with costs
\[
\begin{array}{c|cc}
&\text{tricks}&\text{plays}\\\hline
\C^2&2&6\\
\C^{20}&4&42
\end{array}.
\]
Thus every complete traversal of the three pieces uses thirteen tricks and
$107$ plays, and twelve traversals give
\[
 12(7+2+4)=156,\qquad 12(59+6+42)=1{,}284.
\]

At the anchor with the reservoir player named $A$, each thirteen-trick
round has winner word
\[
\texttt{AAAABBBABAAAB}.
\]
It has least period thirteen: since thirteen is prime and this is a
nonconstant word, there can be no shorter period. Hence any return of the
full rank state must occur after a multiple of thirteen tricks. The
protected module has twelve distinct successive phase words, and the
complete rank state returns on the twelfth traversal. This proves
primitivity without relying on a long unstructured trace.

The two entry tricks of Theorem~\ref{thm:balanced} deliver the cyclic
state occurring four tricks after the displayed protected-module anchor.
The entry packets are of lengths $19$ and $3$, and both return conditions
are checked directly by \texttt{verify.py}.

\section{A general balanced-loader identity}
The special load-bearing role of the long ordinary spacer has a general
form. Here $R$ is a face rank demanding $r$ ordinary cards.

\begin{lemma}[Two-trick balanced loader]\label{lem:loader}
Suppose $r+2\le s\le26$ and a cyclic boundary has the reservoir player
$S$ leading, with queues
\[
S=U\,\C R\C^r,\quad |S|=s,\qquad
L=T\,\C^{53-2s}R\C^r,
\]
where $|U|=|T|=s-r-2$ and these suffixes fit the inventory.
Then the balanced starting queues
\[
\boxed{
S_0=\C^{27+r-s}RU,\qquad
L_0=\C^{26-s}R\C^{r+1}T
}
\]
with $S_0$ leading, reach the given boundary in two completed tricks. The
first has $54+r-2s$ plays, the second $r+2$ plays.
\end{lemma}
\begin{proof}
Both displayed initial hands have length $26$. Put $x=26-s$.
The two leading ordinary prefixes have lengths $x+r+1$ and $x$.
They are played alternately until $L$ exhausts its initial $x$
ordinary cards. Player $S$ plays one further ordinary card, then $L$
plays $R$; $S$ pays $r$ ordinary cards. Thus $L$ takes the
packet $\C^{2x+1}R\C^r$, containing $2x+r+2=54+r-2s$ cards.
Player $L$ now leads $\C$, and $S$ answers with $R$, which $L$
pays by $r$ ordinary cards. Player $S$ takes the packet
$\C R\C^r$ of length $r+2$. The unplayed suffix of $S_0$ is $U$;
the unplayed suffix of $L_0$ is $T$, and its first collection has been
appended after it. The resulting queues are therefore exactly
$U\C R\C^r$ and $T\C^{2x+1}R\C^r$, with $S$ leading.
\end{proof}

For $R=\J$ and $s=18$, the two packets have lengths $19,3$, as in the
balanced example. The lemma is conditional: the required cyclic boundary
must first be found; a large spacer alone does not establish accessibility.

\section{Complete inverse components}
A cycle can have no balanced ancestor even if its periodic hand sizes are
near balance. Given a boundary state after a completed trick, its preceding
pile must be a suffix of the winner's queue. For every nonempty suffix
and either possible previous leader, the challenge rules determine the
sequence of players contributing its cards. Reinsert those contributions
at the two queue fronts and accept the candidate precisely if one forward
trick returns to the chosen boundary.

\begin{proposition}[Computer-certified reverse components]
Under the stated rank-state convention, the complete backward components
of the three cycles have the following sizes:
\begin{center}
\begin{tabular}{@{}lrrrr@{}}
\toprule
Cycle & States & Reconstructed edges & Balanced states & Maximum transient\\
\midrule
$156/1{,}284$ & $7{,}992$ & $7{,}992$ & $72$ & $5$\\
$200/1{,}264$ & $6{,}352$ & $6{,}352$ & $0$ & --\\
$176/1{,}228$ & $6{,}760$ & $6{,}760$ & $0$ & $6$\\
\bottomrule
\end{tabular}
\end{center}
Every balanced state in the first component enters after exactly two
tricks. Neither unrestricted orbit can be reached from a balanced state.
\end{proposition}

The inverse procedure is exhaustive because every possible last trick is
characterized by its collected suffix and previous leader. The verifier
traverses all reconstructed predecessors until its work queue empties,
checks all candidate edges by forward replay, and rejects incomplete
payment or terminal predecessors. The stated sizes are thus reproducible
exhaustive computer-assisted claims about these \emph{particular} components;
they do not classify the components of all cycles. The absence of a balanced
ancestor for the two unrestricted cycles is stronger than merely inspecting
the hand sizes along their periodic orbits.

\section{Scope, verification, and further questions}
The constructions distinguish two different optimization problems:
maximizing the eventual primitive period accessible from a balanced deal,
and maximizing the primitive period over arbitrary nonempty splits. The
examples provide lower bounds $156$ and $200$ respectively; their play
periods have the opposite order. No global maximality is known here.

The local certificates require only Python's standard library. From the
repository root, run
\begin{verbatim}
python3 long_cycles/verify.py \
  --json long_cycles/verification/long_cycle_checks.json
\end{verbatim}
The script verifies the three inventories, both direct simulators, every
primitive state, all protected passages of the balanced module, and all
backward components. The committed JSON is a compact report, not a
replacement for regenerating the checks. The earlier six-trick results and
the present work use the same chronological-pickup convention and rank
state definition \cite{longshort}.

Further targets include balanced-accessible cycles longer than $156$, and
modules whose first-return map carries new coprime face- or ordinary-gap
periods. Changing the reservoir rank or allowing the reservoir state to
switch between passages offers other possibilities. The complete state
space is finite, but the global longest primitive period remains open.

\begin{thebibliography}{9}
\bibitem{paulhus}
M. M. Paulhus, \emph{Beggar My Neighbour},
American Mathematical Monthly \textbf{106} (1999), 162--165.
\href{https://doi.org/10.2307/2589054}{doi:10.2307/2589054}.
\bibitem{casella}
B. Casella, P. M. Anderson, M. Kleber, R. P. Mann, R. Nessler,
W. Rucklidge, S. G. Williams, and N. Wu,
\emph{A Non-Terminating Game of Beggar-My-Neighbor},
American Mathematical Monthly \textbf{132} (2025), 978--994.
\href{https://doi.org/10.1080/00029890.2025.2558484}{doi:10.1080/00029890.2025.2558484}.
\bibitem{andorno}
N. Andorno, G. Cernoia, S. Duiz, and A. Michelangeli,
\emph{The Dynamical Landscape of Beggar-My-Neighbour: Ultra-long Matches,
Loops, and Infinite Matches}, arXiv:2602.23406 (2026).
\href{https://arxiv.org/abs/2602.23406}{arxiv:2602.23406}.
\bibitem{longshort}
C. D. Long, \emph{The Unique Shortest Rank-State Cycle in
Beggar-My-Neighbour: Six Tricks and 38 Individual Plays},
companion manuscript, 2026. \url{https://github.com/long-mathematics/beggar-my-neighbour-period}.
\end{thebibliography}
\end{document}
